\section{Modal Analysis}

\subsection{Analysis Setup}

The Modal module in ANSYS Mechanical has been used to obtain the natural frequencies and mode shapes of the geometry. A minimum of six modes have been extracted.

% Clearly indicate:
% - Type of analysis used
% - Number of calculated modes
% - Final selected mesh
% - Applied boundary conditions
% - Assigned material
% - Main assumptions of the analysis

\subsection{First Six Vibration Modes}

The first six vibration modes are presented below. For each mode, the natural frequency, the observed mode shape and the physical interpretation of the motion are indicated.

\begin{table}[H]
\centering
\caption{Summary of the first six vibration modes.}
\label{tab:modes}
\begin{tabular}{ccc}
\toprule
Mode & Frequency [\si{\hertz}] & Physical description of the mode \\
\midrule
1 & & \\
2 & & \\
3 & & \\
4 & & \\
5 & & \\
6 & & \\
\bottomrule
\end{tabular}
\end{table}

% TODO: Include a screenshot of each mode (6 figures)
% \begin{figure}[H]
%     \centering
%     \includegraphics[width=0.7\textwidth]{mode_1}
%     \caption{Mode 1: description of the mode.}
%     \label{fig:mode_1}
% \end{figure}

\subsection{Interpretation of the Modes}

% For each mode, answer the following questions:
% - Is the mode dominated by bending, torsion, tension-compression or another type of deformation?
% - What is the main direction of motion?
% - Which areas show the largest relative deformation?
% - Does the mode affect the entire geometry or only a local area?
% - Could the mode be relevant if it coincides with an excitation frequency?
% - Are there areas where high stresses could appear in a subsequent dynamic analysis?